% Generated draft: requires mathematical and source review.
\documentclass[11pt]{article}
\usepackage[margin=25mm]{geometry}
\usepackage{fontspec}
\setmainfont{Latin Modern Roman}
\usepackage{amsmath,amssymb,mathtools,unicode-math}
\setmathfont{Latin Modern Math}
\usepackage{enumitem,needspace,xurl,hyperref}
\setlength{\parindent}{0pt}
\setlength{\parskip}{5pt}
\setlength{\emergencystretch}{3em}
\newcommand{\N}{\mathbb{N}}
\newcommand{\Z}{\mathbb{Z}}
\newcommand{\Q}{\mathbb{Q}}
\newcommand{\R}{\mathbb{R}}
\newcommand{\C}{\mathbb{C}}
\newcommand{\F}{\mathbb{F}}
\newcommand{\A}{\mathbb{A}}
\newcommand{\PP}{\mathbb{P}}
\newcommand{\E}{\mathbb{E}}
\newcommand{\eps}{\varepsilon}
\newcommand{\dd}{\,\mathrm{d}}
\begin{document}
\section*{UnsolvedMath ID: 800267}
\textbf{Source problem code:} catalogue entry
\subsection{Definitions and mathematical statement}
For each \(c\in\C\), let
\[
F_c:\C^2\longrightarrow\C^2,\qquad F_c(z,w)=(z^2+c-w,z).
\]
The source specifies that \(F_c\) is a symplectic automorphism: it is biholomorphic and satisfies
\[
F_c^*(\dd z\wedge\dd w)=\dd z\wedge\dd w.
\]
For an integer \(k\geq1\), a periodic orbit is represented by points \(p_0,\ldots,p_k\in\C^2\) satisfying
\[
p_{i+1}=F_c(p_i)\quad(0\leq i<k),\qquad p_k=p_0.
\]
Here \(k\) need not be the least period; the source does not impose a least-period requirement.

\textbf{Question attributed to M. Herman.} Do there exist \(c\in\C\), an integer \(k\geq1\), and such a periodic orbit for which \(p_0\) belongs to a Siegel domain?

\textbf{Unresolved qualification.} The supplied record does not define ``Siegel domain'' in this dynamical setting. In particular, it does not specify the required conjugacy or rotation conditions, or whether the domain is associated with \(F_c\) itself or with an iterate. The question above retains this source terminology; a unique self-contained mathematical formulation requires clarification of that term.
\subsection{Short English statement}
Can some member of the family \(F_c(z,w)=(z^2+c-w,z)\), with \(c\in\C\), have a periodic point lying in a Siegel domain? The intended meaning of ``Siegel domain'' remains unspecified in the supplied record.
\subsection{Sources}
\begin{itemize}
\item Local UnsolvedMath registry, permanent ID \(800267\), source code \texttt{catalogue entry}, title ``Several Complex Variables Problem List --- Problem 66.'' Upstream revision: \texttt{d371ee917ac9647fbb06cd6fc529deec6f3dda82}.
\item Buzzard, Fornaess, and Krantz, \emph{The Several Complex Variables Problem List} (1995), Problem 66. Supplied source URL: \url{https://arxiv.org/abs/math/9509202}. The source question is signed M. Herman; the registry's \texttt{proposed\_by} field names Buzzard, Fornaess, and Krantz.
\item The supplied provenance records access on 2026-10-16 and extraction by \texttt{arxiv\_latex}. This is supplied provenance, not an independent source visit or verification. The registry labels the problem ``open''; that label is supplied curation, not an independently verified mathematical or literature assessment.
\end{itemize}
\end{document}
